\section{Functional Workflows (Flowcharts)}

\subsection{Server dispatch}

\begin{figure}[H]
\centering
\begin{tikzpicture}[
  node distance=6mm,
  every node/.style={font=\small},
  proc/.style={draw,rounded corners,align=center,minimum width=47mm,minimum height=7mm},
  decision/.style={draw,diamond,aspect=2,align=center},
  arr/.style={-{Latex},thick}
]
\node[proc] (start) {Bind TCP socket and listen};
\node[proc,below=of start] (accept) {accept client; allocate session ID};
\node[proc,below=of accept] (thread) {Register Session; start daemon thread};
\node[proc,below=of thread] (recv) {Send 220; receive buffered command};
\node[decision,below=of recv] (auth) {Allowed in current\\authentication state?};
\node[proc,below left=8mm and 12mm of auth] (deny) {Reply 530/503};
\node[proc,below right=8mm and 12mm of auth] (dispatch) {Dispatch command handler};
\node[decision,below=16mm of auth] (end) {QUIT, error,\\or disconnect?};
\node[proc,below=of end] (clean) {Cancel transfer; close sockets; deregister};
\draw[arr] (start)--(accept);
\draw[arr] (accept)--(thread);
\draw[arr] (thread)--(recv);
\draw[arr] (recv)--(auth);
\draw[arr] (auth)--node[above left]{no}(deny);
\draw[arr] (auth)--node[above right]{yes}(dispatch);
\draw[arr] (deny)--(end);
\draw[arr] (dispatch)--(end);
\draw[arr] (end.east) -- ++(12mm,0) |- node[pos=.25,right]{no} (recv.east);
\draw[arr] (end)--node[right]{yes}(clean);
\end{tikzpicture}
\caption{TCP accept, authentication gate, and command dispatch.}
\end{figure}

\subsection{Go-Back-N sender}

\begin{figure}[H]
\centering
\begin{tikzpicture}[
 node distance=6mm,
 every node/.style={font=\small},
 proc/.style={draw,rounded corners,align=center,minimum width=55mm,minimum height=7mm},
 decision/.style={draw,diamond,aspect=2,align=center},
 arr/.style={-{Latex},thick}
]
\node[proc] (prep) {Fragment payload; build DATA packets; FIN on last};
\node[proc,below=of prep] (send) {Fill window: base to base + floor(cwnd)};
\node[decision,below=of send] (event) {ACK, NAK, or RTO?};
\node[proc,below left=9mm and 18mm of event] (ack) {ACK: advance base; update RTO if unambiguous};
\node[proc,below right=9mm and 18mm of event] (loss) {Loss: halve threshold/window; timeout resets cwnd to 1};
\node[proc,below=22mm of event] (update) {ACK: Slow Start or additive increase\\Loss: retransmit outstanding flight};
\node[decision,below=of update] (done) {All packets ACKed?};
\node[proc,below=of done] (ok) {Return success};
\draw[arr] (prep)--(send);
\draw[arr] (send)--(event);
\draw[arr] (event)--node[above left]{ACK}(ack);
\draw[arr] (event)--node[above right]{NAK/RTO}(loss);
\draw[arr] (ack)--(update);
\draw[arr] (loss)--(update);
\draw[arr] (update)--(done);
\draw[arr] (done.east)--++(13mm,0)|-node[pos=.2,right]{no}(send.east);
\draw[arr] (done)--node[right]{yes}(ok);
\end{tikzpicture}
\caption{Cumulative-ACK Go-Back-N sender with congestion response.}
\end{figure}

\subsection{Go-Back-N receiver}

The receiver validates header length and CRC, filters the transfer ID and DAT
flag, accepts only \code{seq == expected_seq}, appends that payload, and
returns a cumulative ACK. Corruption produces NAK. Duplicate/future packets
are discarded and the highest contiguous ACK is repeated. FIN on an accepted
packet terminates reassembly.

\subsection{Data-mode toggle and transfer}

PASV binds at the server and uses client-to-server RDY; PORT binds at the
client and uses server-to-client RDY. Both converge on the same
\code{RDTSender}/\code{RDTReceiver} workflow, then release the per-transfer
UDP socket and return \code{226} or \code{426}.

\subsection{Limitations}

Authentication is basic and has no TLS/password hashing. CRC-32 detects
accidental corruption but is not an HMAC. MODE B is accepted but does not
introduce separate RFC block framing. Hash verification is available but not
forced after every transfer. Payloads are assembled in memory. There is no
per-filename write lock. The Slow Start/AIMD design is educational rather
than a complete TCP congestion-control implementation.